\documentclass[a4paper]{article}
\usepackage{graphicx}
\usepackage{amsmath,amssymb}
\usepackage{hyperref}
\usepackage{minted}
\usepackage{multirow}
\usepackage{float}
\usepackage{todonotes}
\usepackage{forest}

\title{}
\date{}

\begin{document}
    \maketitle
    
    Attention by itself does not know anything about the order of elements in the sequence. The attention operation treats the input as a set rather than an ordered sequence. Therefore the model needs an additional mechanism to represent the position of each element.
    
    This is the role of \textbf{positional encoding}.
    
    For each position \(i\), a positional vector \(p_i\) is created. This vector contains information about the index of the element within the sequence.
    
    The positional vector is added to the embedding:
    
    \[
    z_i = e_i + p_i
    \]
    
    The resulting vector \(z_i\) contains both
    
    \[
    \text{content information} \quad (e_i)
    \]
    
    and
    
    \[
    \text{position information} \quad (p_i)
    \]
    
    One common positional encoding uses sine and cosine functions:
    
    \[
    p_{i,2k} = \sin\!\left(\frac{i}{10000^{2k/d}}\right)
    \]
    
    \[
    p_{i,2k+1} = \cos\!\left(\frac{i}{10000^{2k/d}}\right)
    \]
    
    where \(d\) is the embedding dimension.
    
    After positional encoding is added, the vectors \(z_i\) are used to compute
    
    \[
    Q = W_q z_i, \quad K = W_k z_i, \quad V = W_v z_i
    \]
    
    and the attention mechanism operates on these representations.
    
    In summary, the roles are the following.
    
    Encoding converts raw input values into vectors suitable for neural computation.
    
    Positional encoding injects information about the order of elements in the sequence.
    
    Together they ensure that the Transformer can both interpret the content of each element and understand its position within the sequence.
    
    % \bibliographystyle{plain}
    % \bibliography{/home/donkarlo/Dropbox/repo/shared_research/bibliography/bibliography.bib}
\end{document}